\section{Introduction}
\label{sec:intro}

In 1997 Certicom published a list of elliptic-curve discrete
logarithm challenges~\cite{certicom1997}. Three of them, ECC2K-95,
ECC2K-108 and ECC2K-130, use the Koblitz curve $y^2+xy=x^3+1$ over
$\FF_{2^{97}}$, $\FF_{2^{109}}$ and $\FF_{2^{131}}$ and ask for the $k$
with $Q=[k]P$. Harley's group solved ECC2K-95 in 1998~\cite{harley1998}
and ECC2K-108 in 2000. ECC2K-130 is still open. In 2009 a consortium of
twenty authors published the arithmetic, the iteration function and
per-platform throughput for an attack on
it~\cite{bailey2009,bailey2009sharcs,bos2010cell,fan2010fpga}, and
estimated the work at $2^{60.9}$ iterations. That attack was never
completed.

This paper describes a second attempt, built from scratch on hardware
that did not exist in 2009. We tried to hold it to the standard the
first attempt's reports set, and in places to a stricter one: every
rate is paired against an adjacent control on the same machine, every
target was written down before we tuned toward it, and every claim
has a footnote naming the file in the repository that backs it. None
of this improves on the generic-group bound, and we do not pretend
otherwise. What we can report is the cost of one step, measured on
GPUs, CPUs and an FPGA; the number of steps, measured on small curves
of the same family and extrapolated; and a campaign that has walked
$2^{56.7}$ of the $2^{60.9}$ iterations, most of it before its cloud
fleet was taken away.

Along the way it became clear that most of what made the GPU client
fast was not specific to binary fields. So we also carried the same
optimisations over to Pollard rho on NIST P-256 and P-384, on the same
card, to see which of them transfer.

\paragraph{Contributions.}
\begin{itemize}[leftmargin=1.2em,itemsep=0.3em]
\item A code generator that builds every field routine from one
  symmetric-vector model of $\FF_{2^{131}}$ and checks each routine
  before emitting it. On top of it sit a bitsliced backend at every word
  width from 64 to 512 lanes, a packed CUDA backend that uses the GPU's
  native carry-less products, and a VHDL engine, all of which write
  bit-for-bit compatible records (Sections~\ref{sec:curve} and
  \ref{sec:arith}).
\item Three iteration functions on the $\langle\sigma,-1\rangle$
  classes, with walk constants measured on small Koblitz curves and
  extrapolated. The campaign's walk needs $2^{60.91}$--$2^{60.92}$
  iterations. The table walk's fruitless cycles are measured, priced
  and repaired (Section~\ref{sec:walk}).
\item A GPU ledger of sixteen paired optimisation steps, taking one
  RTX~PRO~6000 from $0.85$ to $14.64$~billion iterations/s for the campaign build, with
  later builds of the same walk at $15.44$ and $15.88$ and a table walk
  at $20.08$~billion iterations/s (a rate whose cost per solve is not yet measured). We
  add a twelve-GPU survey, a roofline, a floor derived from the
  carry-less unit before tuning that rules out $28$~billion iterations/s on one card,
  and every candidate we built and rejected, with its ratio
  (Section~\ref{sec:gpu}).
\item The same optimisation ladder applied to Pollard rho on NIST P-256
  and P-384 on the same RTX~PRO~6000: $1.52\times$ and $2.61\times$ over the
  first working build, to $1.632$ and $0.548$~billion iterations/s, with an account of
  which levers carried over and how their sizes changed
  (Section~\ref{sec:nist}).
\item An inventory of the other GPU arithmetic in the repositories,
  prime-field (secp256k1 rho with order-six folding, baby-step
  giant-step, kangaroo) and binary-field (the ECC2K-95 kernel, pair
  tables, polynomial-system solvers), with their instruction ledgers,
  emulated walk constants and an honest statement of which have run on
  a device (Section~\ref{sec:othergpu}).
\item A pipelined-Karatsuba FPGA engine measured at $8.25$~billion iterations/s on
  one AWS F2, inside a pre-board estimate of $5$--$12$
  (Section~\ref{sec:fpga}).
\item The campaign itself: its configuration, the calibration of the
  distinguished-point interval (including the retraction of a wrong
  one), the storage and replay protocol, fleet economics, and its state
  at $486$~M points and $2^{56.68}$ iterations
  (Section~\ref{sec:campaign}).
\item Structural alternatives priced in the same unit, none of which
  come close: index calculus on the Koblitz curve at $10^7$ times the
  cost of rho, and point decomposition $35$--$51$ bits short of its
  budget (Section~\ref{sec:alternatives}).
\end{itemize}

\paragraph{How we report numbers.} One rule from the repository
governs every section, so we state it once here. A change is measured
as a ratio to a boundary that was written down before the change was
attempted. That boundary is either a \emph{floor}, a counting,
generic-group or hardware argument the method cannot cross, or a
\emph{reference}, the best existing method on the same instance in the
same unit. Every cost the method incurs goes inside the ratio:
precomputation, verification, the inversion a batch shares, and the
work still in flight when a collision lands. A number that prices only
one phase is a stage diagnostic, never a speed. We then classify each
change by what moved. An \emph{advance} lowers the ratio to the floor.
\emph{Engineering} lowers the cost while the ratio to the floor stays
flat. \emph{Relabelling} lowers a headline while the total cost rises.
\emph{Accounting} changes a number without changing the algorithm.
Every row in this paper falls into one of the last three. The same
rule, applied across Pollard rho, baby-step giant-step, kangaroo and
prime-field index calculus, is the subject of a companion draft in the
sibling library~\cite{ecdlppaper}, though this paper does not depend
on it. The mistake we tried hardest to avoid is quoting a rate without
its paired control, its hardware and its launch geometry, and
Sections~\ref{sec:ledger} and \ref{sec:campaign-interval} each record
a case where an earlier draft in the repository made exactly that
mistake.

\paragraph{Availability.} Every number in the paper has a footnote
naming the file it came from. The campaign lives in a public
repository, \thisrepo\ (\url{https://github.com/aburan28/crypto}), and
its files are cited by path from the repository root. A public sibling
library, \sibling\ (\url{https://github.com/aburan28/cryptanalysis}),
holds an independent C golden model of the same field, a digit-serial
FPGA core without vendor timing, an import of the $20$~billion iterations/s GPU client
with its own host re-walk, and the boundary ledger for index calculus;
its paths carry the prefix \texttt{cryptanalysis:}. Two evidence
records cited in Section~\ref{sec:alternatives} are in a third public
repository, the research ledger
(\url{https://github.com/aburan28/crypto-autoresearcher}), with the
prefix \texttt{autoresearcher:}. Live campaign aggregates are published
at \url{https://aburan28.github.io/crypto/status/}, and the snapshot
this paper reports is committed next to the paper source.

\paragraph{What we do not claim.} We have not recovered a logarithm on
the challenge curve, and we do not expect one from the corpus
collected so far; collection is currently paused. We quote no rate for
hardware we did not measure. The L40S collection rate, the RTX~PRO~4500
collection rate under the bridge map and the G7e instance family are
unmeasured, and we say so where they come up. Projections, such as the
$7.2$--$7.4$~billion iterations/s fused build on the 4500, the table walk's
$0.81$--$0.86$ relative cost per solve, and the 136- and 144-engine F2
images, are labelled as projections wherever they appear. The NIST
rates in Section~\ref{sec:nist} are throughput on a walk that will not
finish, and say nothing about the security of those curves.
